% META: {"id":"1SPE-DERGLOBAL-EX-027","chapitre":"1SPE-DERIVATION-GLOBAL","type_objet":"exercice","capacites":["1SPE-DERIVATION-GLOBAL-C3"],"capacites_codes":["C3"],"methodes":["M3"],"parcours":2,"competences":["calculer","modeliser"],"duree_min":12,"mode_creation":"ex_nihilo","sources_inspiration":[],"fichier_tex":"chapitres/1SPE-DERIVATION-GLOBAL/exercices/1SPE-DERGLOBAL-EX-027.tex","corrige_tex":"chapitres/1SPE-DERIVATION-GLOBAL/corriges/1SPE-DERGLOBAL-CO-027.tex","status":"generated"}
% BEGIN-VERIFY
% from sympy import symbols, diff, solve
% t = symbols('t')
% h = -5*t**2 + 20*t + 1
% hp = diff(h, t)
% assert hp == -10*t + 20
% racine = solve(hp, t)
% assert racine == [2]
% assert hp.subs(t, 1) > 0
% assert hp.subs(t, 3) < 0
% assert h.subs(t, 2) == 21
% assert h.subs(t, 0) == 1
% END-VERIFY
\begin{exercice}{1SPE-DERGLOBAL-EX-027}{2}{12}
Un objet est lancé verticalement vers le haut. Sa hauteur (en mètres) au bout de $t$ secondes ($t \in [0\,;\,4]$) est modélisée par :
\[
h(t) = -5t^2 + 20t + 1
\]

\begin{enumerate}
\item Quelle est la hauteur initiale de l'objet ($t = 0$) ?
\item Calculer $h'(t)$ et résoudre $h'(t) = 0$ sur $[0\,;\,4]$.
\item Dresser le tableau de variations de $h$ sur $[0\,;\,4]$.
\item À quel instant l'objet atteint-il sa hauteur maximale ? Quelle est cette hauteur ?
\end{enumerate}
\end{exercice}
